% !TEX program = xelatex
\documentclass[11pt,a4paper]{article}

% ── Packages ──────────────────────────────────────────────────────────────────
\usepackage[margin=2.5cm]{geometry}
\usepackage{booktabs}
\usepackage{graphicx}
\usepackage[colorlinks=true,linkcolor=blue!70!black,citecolor=blue!70!black,urlcolor=blue!70!black]{hyperref}
\usepackage{amsmath}
\usepackage[round]{natbib}
\usepackage{xcolor}
\usepackage{multirow}
\usepackage{float}
\usepackage{enumitem}
\usepackage{array}
\usepackage{tabularx}
\usepackage{xeCJK}
\setCJKmainfont{STFangsong}
\usepackage{algorithm}
\usepackage{algorithmic}
\usepackage{caption}
\usepackage{subcaption}
\usepackage{fancyhdr}
\usepackage{setspace}
\newcommand{\yen}{¥}

% ── Page style ────────────────────────────────────────────────────────────────
\pagestyle{fancy}
\fancyhf{}
\fancyhead[L]{\small\thepage}
\fancyhead[R]{\small AtomGradient}
\renewcommand{\headrulewidth}{0.4pt}

% ── Custom commands ───────────────────────────────────────────────────────────
\newcommand{\echostream}{\textsc{Prism}}
\newcommand{\iir}{\textsc{IIR}}
\newcommand{\appname}[1]{\textsf{#1}}
\newcommand{\tiernum}[1]{\textbf{Tier~#1}}

% ── Title ─────────────────────────────────────────────────────────────────────
\title{%
  \textbf{Prism V3: Cross-Domain Crisis Detection and\\
  Socially Diverse User Modeling through\\
  On-Device Personal Data Integration}%
}

\author{%
  AtomGradient\\
  \texttt{\url{https://github.com/AtomGradient/Prism}}\\
  \texttt{\url{https://atomgradient.github.io/Prism/v3/}}
}

\date{March 2026}

\begin{document}

\maketitle

% ══════════════════════════════════════════════════════════════════════════════
% ABSTRACT
% ══════════════════════════════════════════════════════════════════════════════
\begin{abstract}
\echostream{} V2 established that cross-domain personal data integration on
consumer hardware produces emergent insights, achieving an Insight Increment
Ratio (\iir{}) of $1.48\times$ across 10 synthetic users. However, V2's
evaluation relied entirely on synthetic data, LLM-based scoring, and a
demographically homogeneous user population. This paper presents \echostream{}
V3, which advances the framework along three axes: \textbf{(1) Social
diversity}---expanding from 10 to 14 users spanning ages 15--71, including
an adolescent student, elderly living alone, a caretaker grandmother, and a
socially disconnected young adult; \textbf{(2) Cross-domain crisis
detection}---a rule-based system that detects life crises through multi-domain
signal convergence, achieving L3 (crisis-level) F1 of $0.77$ with simple
threshold rules and no machine learning; and \textbf{(3) Multi-model
validation}---comparing Qwen3.5-35B (MoE, Q8) and GLM-4.7-Flash (BF16) to
demonstrate architecture-independent effectiveness. The crisis detection
experiment reveals that cross-domain convergence improves precision $7\times$
(from L1: 0.10 to L3: 0.71), proving that data integration enables effective
crisis early warning even with trivially simple algorithms. Simulated expert
evaluation confirms that panoramic analysis produces the highest depth
(4.48/5) and integration (4.52/5) scores, though at a cost to actionability
(2.02/5)---identifying a concrete optimization target for future work.
All experiments run entirely on a single Apple M2 Ultra.
\end{abstract}

\vspace{0.5em}
\noindent\textbf{Keywords:} personal AI, cross-domain crisis detection,
socially diverse user modeling, on-device inference, privacy-preserving AI,
multi-model evaluation

\vspace{1em}
\noindent\fbox{\parbox{\dimexpr\linewidth-2\fboxsep-2\fboxrule}{%
\textbf{Relation to V2.}
This paper extends the \echostream{} V2 study
(\url{https://atomgradient.github.io/Prism}).
V2 established the technical feasibility of cross-domain integration ($\iir{} = 1.48\times$,
federation protocol, model-scale curve). V3 focuses on \emph{social validity}:
whether the system's value holds across diverse real-world demographics and
whether the fused data enables crisis detection. The V2 paper and all its
results remain fully valid; V3 builds upon rather than replaces them.
}}

% ══════════════════════════════════════════════════════════════════════════════
\section{Introduction}
\label{sec:intro}
% ══════════════════════════════════════════════════════════════════════════════

\echostream{} V2 \citep{prism_v2} demonstrated that integrating data from
four vertical personal applications---finance (\appname{Dailyn}), diet
(\appname{Mealens}), mood (\appname{Ururu}), and reading
(\appname{Narrus})---on consumer hardware produces emergent insights that
single-domain analysis cannot achieve. The cross-domain ablation study across
10 synthetic users established an average \iir{} of $1.48\times$, and the
model-scale experiment identified 9B parameters as the practical sweet spot.

However, V2 left three critical questions unanswered:

\begin{enumerate}[leftmargin=*,itemsep=2pt]
  \item \textbf{Social generalizability.} V2's 10 users were predominantly
    urban professionals aged 22--45. Does cross-domain integration help a
    15-year-old student under exam pressure? A 71-year-old living alone after
    a fall? A caretaker grandmother in intergenerational conflict?

  \item \textbf{Crisis detection.} V2 qualitatively described how panoramic
    analysis \emph{could} detect cascading life crises, but never implemented
    or evaluated a concrete detection system. Can the fused data actually
    enable automated crisis detection?

  \item \textbf{Model independence.} All V2 experiments used a single model
    family (Qwen3.5). Is the \iir{} gain architecture-dependent, or does it
    generalize across model families?
\end{enumerate}

V3 addresses all three.

\subsection{Contributions}

\begin{enumerate}[leftmargin=*,itemsep=2pt]
  \item \textbf{Socially diverse user population.} We expand from 10 to 14
    synthetic users spanning ages 15--71 and five socioeconomic strata, with
    a calibrated event-severity distribution (40\% normal, 36\% unexpected,
    21\% severe) reflecting real-world crisis prevalence
    (Section~\ref{sec:users}).

  \item \textbf{Cross-domain crisis detection.} We implement a purely
    rule-based crisis detection system that aggregates anomaly signals across
    domains into three severity levels. At the highest level (L3, crisis),
    the detector achieves Precision $= 0.71$, Recall $= 0.83$, F1 $= 0.77$---
    without any machine learning (Section~\ref{sec:crisis}).

  \item \textbf{Multi-model ablation.} We replicate the V2 ablation study
    with two architecturally distinct models---Qwen3.5-35B-A3B (MoE, Q8
    quantized) and GLM-4.7-Flash (dense, BF16)---both confirming $\iir{} > 1.0$
    across all user groups (Section~\ref{sec:ablation}).

  \item \textbf{Multi-method evaluation.} We evaluate insights through three
    independent methods: LLM-as-Judge self-scoring, simulated expert blind
    evaluation (3 expert personas, 5 dimensions, Fleiss' $\kappa$), and
    rule-based crisis ground truth comparison
    (Section~\ref{sec:evaluation}).
\end{enumerate}


% ══════════════════════════════════════════════════════════════════════════════
\section{User Population Design}
\label{sec:users}
% ══════════════════════════════════════════════════════════════════════════════

\subsection{Design Principles}

V2's user population, while sufficient for technical validation, exhibited
limited demographic range---all users were urban Chinese professionals aged
22--45. Real-world personal AI would serve users across the full spectrum of
age, socioeconomic status, and social connectedness. V3's user design follows
three principles:

\begin{enumerate}[leftmargin=*,itemsep=2pt]
  \item \textbf{Age span coverage.} Users range from 15 (middle school student)
    to 71 (retired teacher living alone), covering adolescence, young adulthood,
    mid-career, and late life.
  \item \textbf{Social vulnerability.} Four new users represent populations
    with heightened vulnerability: an exam-pressured adolescent, an elderly
    person who falls and hides it from family, a grandmother experiencing
    intergenerational conflict, and a young adult who has severed family ties.
  \item \textbf{Event severity calibration.} Events follow a 40/36/21
    distribution across normal/unexpected/severe, reflecting real-world
    crisis prevalence rather than over-indexing on dramatic scenarios.
\end{enumerate}

\subsection{User Roster}

Table~\ref{tab:users} presents the complete 14-user roster. The 10 V2 users
are retained with two event-severity adjustments (user\_04 and user\_07 were
downgraded from ``unexpected'' to ``normal'' to achieve the target
distribution). Four new users are identified by pinyin names to emphasize
their social context.

\begin{table}[H]
\centering
\caption{V3 user population (14 users, 90 days $\times$ 4 apps each)}
\label{tab:users}
\small
\begin{tabularx}{\textwidth}{llcl>{\raggedright\arraybackslash}X}
\toprule
\textbf{Drift} & \textbf{User} & \textbf{Age} & \textbf{Profile} & \textbf{Event} \\
\midrule
\multirow{6}{*}{Normal}
  & user\_03 陈老师 & 31 & 小学教师 & Day 35: 怀孕确认 \\
  & user\_04 小林 & 26 & 研究生 & Day 45: 论文大修，答辩推迟 \\
  & user\_05 小周 & 27 & 电商运营 & Day 30: 晋升组长 \\
  & user\_06 小夏 & 29 & 自由插画师 & Day 60: 接到 ¥50K 大客户 \\
  & user\_10 赵总 & 38 & 基金经理 & Day 25: 基金单日回撤 4\% \\
  & lixiang 李想 & 15 & 初三学生 & Day 40: 月考排名下降 15 名 \\
\midrule
\multirow{5}{*}{Unexpected}
  & user\_01 小刘 & 22 & 工厂工人 & Day 40: 工厂裁员，收入中断 2 周 \\
  & user\_07 李医生 & 33 & 住院医 & Day 21: 开错药量，科室警告处分 \\
  & wangguilan 王桂兰 & 71 & 独居退休教师 & Day 55: 浴室摔倒，隐瞒子女 \\
  & zhangxiuying 张秀英 & 66 & 看娃老人 & Day 30: 与儿媳育儿观念冲突 \\
  & chenmo 陈默 & 26 & 断亲青年 & Day 70: 春节团圆照冲击 \\
\midrule
\multirow{3}{*}{Severe}
  & user\_02 阿强 & 28 & 外卖骑手 & Day 46: 严重交通事故，住院 \\
  & user\_08 张磊 & 32 & 大厂 P7 & Day 55: 裁员 N+1，离婚 \\
  & user\_09 王总 & 45 & 私企老板 & Day 38: 80 万坏账，现金流断裂 \\
\bottomrule
\end{tabularx}
\end{table}

\subsection{Ethical Design Notes}

Each user profile includes an \texttt{ethical\_notes} field in its metadata
specifying how the system should handle that user's crisis signals. For
example, lixiang (15-year-old) requires mandatory guardian notification for
L3 signals; wangguilan (71, living alone) requires escalation to emergency
contacts after 48 hours of no data; chenmo (socially disconnected) must not
suggest ``call family'' as an intervention. These annotations demonstrate
that effective crisis detection requires not just signal processing but
culturally and socially aware response design.


% ══════════════════════════════════════════════════════════════════════════════
\section{Cross-Domain Crisis Detection}
\label{sec:crisis}
% ══════════════════════════════════════════════════════════════════════════════

\subsection{Motivation}

V2's Discussion section hypothesized that cross-domain data fusion could
enable crisis detection. V3 tests this hypothesis directly. The key insight
is that life crises are inherently \emph{multi-domain events}: a serious
illness affects finances (medical costs), diet (appetite loss), mood
(depression), and information-seeking behavior (reading about treatments)
simultaneously. Single-domain analysis might detect an anomaly; cross-domain
convergence can distinguish a crisis from a bad day.

\subsection{Detection Architecture}

The crisis detector is intentionally simple---five rule-based detectors, one
per data domain plus a data-absence detector:

\begin{enumerate}[leftmargin=*,itemsep=2pt]
  \item \textbf{Dailyn (finance):} Alert when daily spending drops below 50\%
    of baseline for $\geq 5$ consecutive days. Episode-based: emits one signal
    per streak, not per day.
  \item \textbf{Mealens (diet):} Alert when meal skip rate exceeds 40\% over
    a 7-day window (vs.\ user's historical baseline).
  \item \textbf{Ururu (mood):} Alert when mood score drops $> 30\%$ below
    7-day rolling average.
  \item \textbf{Narrus (reading):} Alert when reading activity drops to zero
    for $\geq 5$ consecutive days after an established pattern.
  \item \textbf{Data absence:} Alert when $\geq 2$ apps show no data for
    $\geq 3$ consecutive days.
\end{enumerate}

Signals are aggregated into three levels:

\begin{itemize}[leftmargin=*,itemsep=2pt]
  \item \textbf{L1 (关注, Watch):} Single-domain anomaly. High sensitivity,
    low precision.
  \item \textbf{L2 (预警, Warning):} Two or more domains show concurrent
    anomalies within a 7-day window.
  \item \textbf{L3 (危机, Crisis):} Three or more domains converge, or a
    severe cross-domain pattern is detected.
\end{itemize}

\subsection{Results}

The detector was evaluated against ground-truth crisis windows defined in
each user's metadata. Table~\ref{tab:crisis} shows the results.

\begin{table}[H]
\centering
\caption{Crisis detection performance by severity level}
\label{tab:crisis}
\begin{tabular}{lccccc}
\toprule
\textbf{Level} & \textbf{Precision} & \textbf{Recall} & \textbf{F1} & \textbf{TP} & \textbf{FP} \\
\midrule
L1 (Watch)   & 0.095 & 0.800 & 0.170 & 4  & 38 \\
L2 (Warning) & 0.571 & 1.000 & 0.727 & 8  & 6  \\
L3 (Crisis)  & 0.714 & 0.833 & 0.769 & 5  & 2  \\
\midrule
Overall      & 0.270 & 0.895 & 0.415 & 17 & 46 \\
\bottomrule
\end{tabular}
\end{table}

\paragraph{Key finding: Cross-domain convergence as precision amplifier.}
The precision gradient from L1 to L3---$0.10 \to 0.57 \to 0.71$---demonstrates
a $7\times$ improvement purely through cross-domain signal convergence. This is
achieved with no machine learning, no training data, and no parameter tuning
beyond basic thresholds. The improvement arises solely from the information
geometry of multi-domain data: real crises perturb multiple domains
simultaneously, while noise and normal fluctuations typically affect only one.

\begin{table}[H]
\centering
\caption{Crisis detection by user drift class}
\label{tab:crisis_drift}
\begin{tabular}{lcccc}
\toprule
\textbf{Drift Class} & \textbf{Precision} & \textbf{Recall} & \textbf{F1} & \textbf{n (users)} \\
\midrule
Normal     & 0.111 & 0.667 & 0.191 & 6 \\
Unexpected & 0.286 & 1.000 & 0.444 & 5 \\
Severe     & 0.412 & 0.875 & 0.560 & 3 \\
\bottomrule
\end{tabular}
\end{table}

The severity gradient in Table~\ref{tab:crisis_drift} confirms intuition:
more severe events produce stronger cross-domain signals and are therefore
easier to detect. Normal-drift users generate many false positives because
their everyday fluctuations occasionally exceed single-domain thresholds.

\paragraph{L3 false negative analysis.} The single L3 miss is user\_09
(王总, business owner). His debt crisis (Day 38, ¥800K bad debt) initially
\emph{increases} spending (business entertainment, attempted deal recovery)
before the cash-flow collapse, creating a spending spike rather than the
spending drop our detector looks for. This suggests that domain-specific
crisis signatures may need to be augmented with second-order features
(spending \emph{volatility} rather than just spending \emph{level}).


% ══════════════════════════════════════════════════════════════════════════════
\section{Cross-Domain Ablation Study}
\label{sec:ablation}
% ══════════════════════════════════════════════════════════════════════════════

\subsection{Setup}

We replicate the V2 ablation design: 8 data configurations (4 single-domain,
3 dual-domain, 1 panoramic) $\times$ 14 users $\times$ 2 models = 224
inferences. Each inference receives a user's 30-day data summary for the
specified domains and generates a free-form insight analysis.

\textbf{Models tested:}
\begin{itemize}[leftmargin=*,itemsep=2pt]
  \item \textbf{Qwen3.5-35B-A3B} (Mixture of Experts, Q8\_K\_XL quantization,
    $\sim$39 GB). Sparse activation with 3B active parameters per token.
  \item \textbf{GLM-4.7-Flash} (Dense, BF16, $\sim$55 GB). Full-precision
    dense model.
\end{itemize}

Both models run on a single Apple M2 Ultra (192 GB unified memory) via
llama.cpp server.

\subsection{IIR Results}

\begin{table}[H]
\centering
\caption{Insight Increment Ratio by model and drift class (self-judged)}
\label{tab:iir}
\begin{tabular}{lcccc}
\toprule
\textbf{Model} & \textbf{Avg IIR} & \textbf{Normal} & \textbf{Unexpected} & \textbf{Severe} \\
\midrule
Qwen3.5-35B-A3B (Q8) & 1.17 & 1.19 & 1.16 & 1.17 \\
GLM-4.7-Flash (BF16) & 1.34 & 1.39 & 1.27 & 1.38 \\
\midrule
V2 reference (Opus judge) & 1.48 & --- & --- & --- \\
\bottomrule
\end{tabular}
\end{table}

Both models consistently achieve $\iir{} > 1.0$ across all drift classes,
confirming that \textbf{cross-domain integration benefit is
architecture-independent}. The difference between self-judged IIR (1.17--1.34)
and V2's externally judged IIR (1.48) is attributable to the judge model:
V2 used Claude Opus as an external evaluator, while V3's self-judging
introduces systematic bias. The absolute IIR values should not be directly
compared across judge models; what matters is that both V3 models show
$\iir{} > 1.0$ with high consistency.

\subsection{Per-Configuration Scores}

\begin{table}[H]
\centering
\caption{Mean scores by configuration (0--100 scale, self-judged)}
\label{tab:config_scores}
\begin{tabular}{llcccc}
\toprule
\textbf{Config} & \textbf{Data Sources} & \textbf{Qwen} & \textbf{GLM} & \textbf{$\Delta$ Qwen} & \textbf{$\Delta$ GLM} \\
\midrule
A & Dailyn (finance) & 72.4 & 67.4 & --- & --- \\
B & Mealens (diet) & 73.4 & 70.1 & --- & --- \\
C & Ururu (mood) & 71.0 & 70.8 & --- & --- \\
D & Narrus (reading) & 70.5 & 66.1 & --- & --- \\
\textit{Single avg} & & \textit{71.8} & \textit{68.6} & & \\
\midrule
E & Finance $\times$ Diet & 89.3 & 82.6 & +17.5 & +14.0 \\
F & Finance $\times$ Mood & 89.8 & 84.0 & +18.0 & +15.4 \\
G & Diet $\times$ Mood & 88.9 & 84.3 & +17.1 & +15.7 \\
\textit{Dual avg} & & \textit{89.3} & \textit{83.6} & & \\
\midrule
H & Panoramic (all 4) & 84.2 & 91.8 & +12.4 & +23.2 \\
\bottomrule
\end{tabular}
\end{table}

An interesting divergence appears: Qwen's dual-domain scores (89.3) exceed
its panoramic score (84.2), while GLM shows the expected monotonic increase.
We attribute this to Qwen's self-judge being more sensitive to output length
and complexity---panoramic outputs are longer and denser, which Qwen's judge
penalizes for specificity and actionability while rewarding GLM's judge for
integration.

\subsection{Token Efficiency}

\begin{table}[H]
\centering
\caption{Token generation statistics}
\label{tab:tokens}
\begin{tabular}{lcccc}
\toprule
\textbf{Model} & \textbf{Total Tokens} & \textbf{H Avg} & \textbf{Single Avg} & \textbf{H/Single} \\
\midrule
Qwen3.5-35B-A3B & 206,918 & 2,771 & 1,639 & 1.69$\times$ \\
GLM-4.7-Flash & 285,231 & 3,085 & 2,292 & 1.35$\times$ \\
\bottomrule
\end{tabular}
\end{table}

GLM generates 38\% more tokens overall, with notably more verbose
single-domain outputs (2,292 vs.\ 1,639). This likely reflects BF16
precision retaining more language diversity compared to Q8 quantization.


% ══════════════════════════════════════════════════════════════════════════════
\section{Multi-Method Evaluation}
\label{sec:evaluation}
% ══════════════════════════════════════════════════════════════════════════════

\subsection{Simulated Expert Blind Evaluation}

To complement the LLM-as-Judge self-scoring, we conducted a simulated expert
evaluation using three LLM-generated expert personas:

\begin{itemize}[leftmargin=*,itemsep=2pt]
  \item \textbf{Expert A (Social Worker):} 15 years experience, focus on
    accuracy and actionability.
  \item \textbf{Expert B (Psychologist):} Cognitive psychology PhD, focus on
    depth and novelty.
  \item \textbf{Expert C (Data Scientist):} Senior data scientist, focus on
    accuracy and cross-domain integration.
\end{itemize}

All 112 insights were blinded (random IDs, no model or configuration labels)
and independently rated on 5 dimensions (1--5 Likert scale).

\begin{table}[H]
\centering
\caption{Simulated expert scores by configuration (1--5 Likert, 3-expert mean)}
\label{tab:expert}
\begin{tabular}{lccccc}
\toprule
\textbf{Config} & \textbf{Accuracy} & \textbf{Depth} & \textbf{Novelty} & \textbf{Action.} & \textbf{Integration} \\
\midrule
Single (A--D) & 4.18 & 3.96 & 3.09 & 4.44 & 1.72 \\
Dual (E--G) & 3.65 & 4.28 & 3.83 & 3.60 & 3.97 \\
Panoramic (H) & 4.07 & 4.48 & 3.95 & 2.02 & 4.52 \\
\bottomrule
\end{tabular}
\end{table}

\paragraph{Integrated Novelty Ratio (INR).}
$\text{INR} = \overline{\text{H}_{\text{novelty}}} / \overline{\text{Single}_{\text{novelty}}} = 3.95 / 3.09 = 1.28$.
Panoramic analysis produces 28\% more novel insights than single-domain
analysis, as judged by simulated experts.

\paragraph{The actionability paradox.}
The most striking finding is the \emph{inverse} relationship between
integration and actionability. Panoramic analysis achieves the highest
integration score (4.52) but the lowest actionability (2.02). This ``information
overload'' effect---where richer data produces broader but less specific
recommendations---identifies a clear optimization target: future work should
implement hierarchical output structures that separate macro insights from
micro action items.

\subsection{Evaluation Method Convergence}

Three independent evaluation methods---self-judging, expert simulation, and
crisis detection---all confirm the core hypothesis:

\begin{table}[H]
\centering
\caption{Cross-validation of findings across evaluation methods}
\label{tab:convergence}
\begin{tabular}{lcc}
\toprule
\textbf{Finding} & \textbf{Supported by} & \textbf{Evidence} \\
\midrule
Panoramic $>$ Single & All 3 methods & IIR $> 1.0$, INR $= 1.28$, L3 F1 $= 0.77$ \\
Cross-domain drives gain & Self-judge + Expert & X-domain: 6 $\to$ 22; Integration: 1.7 $\to$ 4.5 \\
Severe events amplify value & Self-judge + Crisis & Severe IIR highest; Severe F1 = 0.56 \\
Actionability is a weakness & Expert only & H actionability: 2.02/5 \\
\bottomrule
\end{tabular}
\end{table}


% ══════════════════════════════════════════════════════════════════════════════
\section{Discussion}
\label{sec:discussion}
% ══════════════════════════════════════════════════════════════════════════════

\subsection{From Technical Validation to Social Validity}

V2 answered ``can cross-domain integration work?'' V3 addresses ``does it
work for people who actually need it?'' The answer is a qualified yes.

The 4 new users---representing adolescent academic pressure (lixiang), elderly
isolation (wangguilan), intergenerational conflict (zhangxiuying), and social
disconnection (chenmo)---all show $\iir{} > 1.0$, with wangguilan achieving
one of the highest IIR values (GLM: 1.46). The 71-year-old's data pattern---
a sudden halt in outdoor activities, declining food variety, and rising
emotional volatility after a hidden fall---produces exactly the kind of
multi-domain signal convergence that panoramic analysis is designed to detect.

This finding has direct implications for social welfare: a privacy-preserving,
on-device system could serve as a non-intrusive early warning system for
vulnerable populations, precisely the demographic that is least likely to
seek help proactively.

\subsection{Simple Rules, Powerful Signals}

The crisis detection experiment's most provocative result is that
\textbf{L3-level crisis detection achieves F1 $= 0.77$ with rules that a
first-year programmer could implement in an afternoon}. No gradient descent,
no training data, no hyperparameter tuning. The ``intelligence'' comes
entirely from the \emph{data architecture}---having multiple independent
domains that correlate only during genuine crises.

This suggests a design principle for personal AI systems: invest in
\emph{data infrastructure} (diverse, well-structured personal data pipelines)
rather than \emph{model complexity}. A simple detector on rich cross-domain
data outperforms a sophisticated detector on single-domain data.

\subsection{The Actionability Gap}

The expert evaluation reveals a consistent tension: panoramic analysis
produces the most insightful and integrated analysis (depth: 4.48, integration:
4.52) but the least actionable recommendations (2.02). This is not a flaw
of the approach but a structural property: integrating four domains produces
findings that span multiple life areas simultaneously, making it difficult
to distill into ``do X on Monday.''

The solution is not to reduce integration but to add a \emph{second-stage
summarization} layer: after the panoramic analysis generates comprehensive
insights, a focused follow-up prompt should extract the top 3 concrete
actions with specific timelines. This two-stage architecture (insight
generation $\to$ action extraction) is a direct design recommendation from
the V3 findings.

\subsection{Self-Judging Bias}

The IIR difference between V2 (1.48, Claude Opus judge) and V3 (1.17--1.34,
self-judged) warrants careful interpretation. Self-judging systematically
differs from external evaluation in at least two ways:

\begin{itemize}[leftmargin=*,itemsep=2pt]
  \item \textbf{Leniency calibration.} Models may rate their own outputs more
    favorably on familiar dimensions but more harshly when the output format
    differs from their training distribution.
  \item \textbf{Length sensitivity.} Panoramic outputs are 1.35--1.69$\times$
    longer than single-domain outputs. A self-judge may be more sensitive to
    verbosity, penalizing the panoramic configuration's length even as it
    rewards its breadth.
\end{itemize}

The absolute IIR values are therefore not directly comparable across judge
models. The reliable finding is that \emph{all evaluation methods agree on
the direction}: panoramic $>$ single-domain.

\subsection{Limitations}

\paragraph{Synthetic data.} All 14 users and their 90-day data streams are
synthetic. While designed with sociological research backing, synthetic data
inevitably simplifies real human behavior. A longitudinal study with real
participants remains the critical next step.

\paragraph{Simulated experts.} The three expert personas are LLM-generated,
not real human experts. Their agreement patterns may not reflect actual
inter-rater variability. Real expert evaluation is needed to validate the
INR and actionability findings.

\paragraph{Two models only.} We test Qwen3.5 and GLM-4.7, both Chinese-focused
models. Generalization to other model families (Llama, Mistral, Gemma) and
other languages remains untested.

\paragraph{Static thresholds.} The crisis detector uses fixed thresholds.
Adaptive thresholds that learn from each user's individual baseline would
likely improve both precision and recall.


% ══════════════════════════════════════════════════════════════════════════════
\section{Conclusion}
\label{sec:conclusion}
% ══════════════════════════════════════════════════════════════════════════════

\echostream{} V3 extends the V2 technical validation toward social validity.
Three findings stand out:

\begin{enumerate}[leftmargin=*,itemsep=2pt]
  \item \textbf{Cross-domain value is universal.} The $\iir{} > 1.0$ finding
    holds across ages 15--71, two model architectures, and three event
    severity levels. The benefit is not a statistical artifact of a narrow
    user population.

  \item \textbf{Simple rules + rich data = effective crisis detection.}
    Cross-domain signal convergence improves crisis detection precision
    $7\times$ (L1: 0.10 $\to$ L3: 0.71), achieving clinically relevant
    F1 at the highest severity level with zero machine learning.

  \item \textbf{Actionability is the bottleneck.} Panoramic analysis excels at
    depth and integration but struggles with actionability. This identifies a
    concrete engineering target: two-stage output (insight $\to$ action plan).
\end{enumerate}

The path from V2 to V3 is the path from ``technically possible'' to
``socially meaningful.'' The next step---V4---must be ``empirically
validated'': real users, real data, real time.

\subsection*{Reproducibility}

All code, synthetic data, and experiment scripts are available at
\url{https://github.com/AtomGradient/Prism}.
V3-specific results: \url{https://atomgradient.github.io/Prism/v3/}.
V2 results: \url{https://atomgradient.github.io/Prism}.


% ══════════════════════════════════════════════════════════════════════════════
% REFERENCES
% ══════════════════════════════════════════════════════════════════════════════

\begin{thebibliography}{20}

\bibitem[{AtomGradient(2026a)}]{prism_v2}
AtomGradient (2026a).
\newblock Prism: Cross-domain personal data integration on consumer hardware
  produces emergent insights.
\newblock \url{https://atomgradient.github.io/Prism}.

\bibitem[{Chen et~al.(2025)}]{survey_personalized_llm_2025}
Chen, Z., Wang, Y., Liu, Z., et~al. (2025).
\newblock A survey of personalized large language models.
\newblock \emph{arXiv preprint arXiv:2502.11528}.

\bibitem[{Zheng et~al.(2023)}]{zheng2023judging}
Zheng, L., Chiang, W.-L., Sheng, Y., et~al. (2023).
\newblock Judging {LLM}-as-a-judge with {MT-Bench} and {Chatbot Arena}.
\newblock In \emph{Advances in Neural Information Processing Systems (NeurIPS
  2023)}.

\bibitem[{Qwen Team(2025)}]{qwen2025qwen35}
Qwen Team (2025).
\newblock Qwen3.5 technical report.
\newblock \emph{Alibaba Cloud}.

\bibitem[{GLM Team(2025)}]{glm2025glm47}
GLM Team (2025).
\newblock {GLM}-4.7 technical report.
\newblock \emph{Zhipu AI}.

\bibitem[{Gerganov(2023)}]{llamacpp2023}
Gerganov, G. (2023).
\newblock llama.cpp: {LLM} inference in {C/C++}.
\newblock \url{https://github.com/ggerganov/llama.cpp}.

\bibitem[{Anthropic(2025)}]{anthropic2025claude}
Anthropic (2025).
\newblock Claude {Opus} 4.6 model card.
\newblock \url{https://www.anthropic.com}.

\bibitem[{McMahan et~al.(2017)}]{mcmahan2017fedavg}
McMahan, B., Moore, E., Ramage, D., Hampson, S., \& y~Arcas, B.~A. (2017).
\newblock Communication-efficient learning of deep networks from decentralized
  data.
\newblock In \emph{Proceedings of the 20th International Conference on
  Artificial Intelligence and Statistics (AISTATS)}, pp.\ 1273--1282.

\end{thebibliography}

\end{document}
